
Comme \(\chi(G(n))=\sum_{i=0}^{r-1} a_{i+1}\binom{n+i}{i}\) par 1.7, $J$ est $b_1$-r\'egulier par 6.2\,; donc \(H^q(I(p))=0\) pour $q\ge2$ et $p\ge b_1-q$ par 1.4 (i). La suite exacte \(H^q({O}_P^{\oplus M}(p))\longrightarrow H^q(F(m+p))\longrightarrow H^{q+1}(I(p))\) montre alors que $F(m)$ est $(b_1-1)$-r\'egulier\,; d'o\`u (i).

L'assertion (i) implique que \(0\le h^0(F(b))=\chi(F(b))=\sum_{i=0}^{r}a_i\binom{b+i}{i}\)\,; d'o\`u (ii). D'apr\`es 2.10 et 1.9,
\[
   \chi(F(b))=\sum_{i=0}^{r} a_{m,i}\binom{b_1-1+i}{i}\le b_1+\sum_{i=0}^{r} c_{m,i}\binom{b_1-1+i}{i}=b_0\,;
\]
donc, $h^0(F(-1))\le h^0(F(b))\le b_0$. Par r\'ecurrence sur $r$, on peut supposer que $G$ est un $(b_1,\ldots,b_r)$-faisceau\,; d'o\`u (iii) gr\^ace \`a 1.6.2.

\textbf{Lemme 6.5.} \emph{Sous les conditions 6.1, on suppose $F$ sans cycle premier associ\'e de dimension $0$.}
\begin{enumerate}[label=(\roman*)]
\item \emph{Si $h^0(F)\ge1$, alors $h^0(F(-1))\le h^0(F)-1$.}
\item \emph{$H^0(F(-n))=0$ pour $n\ge h^0(F)$.}
\item \emph{Soient une suite exacte $0\to F(-1)\to F\to G\to0$ et un entier $n_0\ge0$ tels que $H^0(G(-n_0))=0$\,; alors $H^0(F(-n_0))=0$ et $h^0(F)\le n_0.h^0(G)$.}
\end{enumerate}
En effet, soit $\sigma\in H^0(F)$, $\sigma\ne0$; alors $\sigma$ d\'efinit un sous-faisceau non nul $F'={O}_P.\sigma$ de $F$, donc $s=\dim.\operatorname{supp}(F')$ est $\ge1$. Prenons un diagramme commutatif \`a lignes et colonnes exactes
\[
\begin{tikzcd}[column sep=large,row sep=large]
&0 \arrow[d]&0 \arrow[d]&0 \arrow[d]&\\
0 \arrow[r]&F'(-1) \arrow[r] \arrow[d]&F' \arrow[r] \arrow[d]&G' \arrow[r] \arrow[d]&0\\
0 \arrow[r]&F(-1) \arrow[r]&F \arrow[r]&G \arrow[r]&0 & .
\end{tikzcd}
\]
